%% manual.tex --- ratvalue: Manual do Usuário
%% Compile: pdflatex manual.tex  (requer pacotes padrão texlive-full)
%% Codificação: UTF-8

\documentclass[12pt, a4paper, oneside]{book}

%% ── Pacotes ────────────────────────────────────────────────────────────────────
\usepackage[utf8]{inputenc}
\usepackage[T1]{fontenc}
\usepackage[brazil]{babel}
\usepackage{amsmath, amssymb, amsthm}
\usepackage{listings}
\usepackage{xcolor}
\usepackage{booktabs}
\usepackage[margin=2.5cm]{geometry}
\usepackage[colorlinks=true, linkcolor=blue!70!black,
            citecolor=green!50!black, urlcolor=blue]{hyperref}
\usepackage{microtype}
\usepackage{enumitem}
\usepackage{tcolorbox}
\usepackage{tabularx}
\usepackage{array}
\usepackage{mathtools}
\usepackage{tikz}
\usepackage{lmodern}

\tcbuselibrary{skins, breakable}

%% ── Estilo de código C++ ───────────────────────────────────────────────────────
\definecolor{cppblue}{rgb}{0.0,0.25,0.6}
\definecolor{cppgreen}{rgb}{0.0,0.45,0.1}
\definecolor{cppgray}{rgb}{0.45,0.45,0.45}
\definecolor{cppbg}{rgb}{0.97,0.97,0.97}
\definecolor{warningbg}{rgb}{1.0,0.97,0.88}

\lstdefinestyle{cpp}{
    language=C++,
    basicstyle=\ttfamily\footnotesize,
    keywordstyle=\color{cppblue}\bfseries,
    commentstyle=\color{cppgreen}\itshape,
    stringstyle=\color{red!70!black},
    directivestyle=\color{purple},
    showstringspaces=false,
    breaklines=true,
    frame=single,
    backgroundcolor=\color{cppbg},
    rulecolor=\color{gray!40},
    tabsize=4,
    xleftmargin=1em,
    xrightmargin=0.5em,
    numbers=left,
    numberstyle=\tiny\color{cppgray},
    numbersep=6pt,
    morekeywords={int64_t, int32_t, uint8_t, __int128, expected, unexpected,
                  string_view, vector, pair, optional},
}

%% ── Caixas especiais ───────────────────────────────────────────────────────────
\newtcolorbox{definicao}[1]{
  colback=blue!5, colframe=blue!50!black,
  title={\textbf{Definição:} #1}, fonttitle=\small\bfseries,
  breakable}

\newtcolorbox{atencao}{
  colback=warningbg, colframe=orange!80!black,
  title={\textbf{Atenção}}, fonttitle=\small\bfseries,
  breakable}

\newtcolorbox{exemplo}[1]{
  colback=green!5, colframe=green!50!black,
  title={\textbf{Exemplo:} #1}, fonttitle=\small\bfseries,
  breakable}

%% ── Comandos matemáticos ───────────────────────────────────────────────────────
%% \ensuremath garante que funcionam tanto em modo texto quanto matemático
\newcommand{\EBIT}{\ensuremath{\mathrm{EBIT}}}
\newcommand{\FCFF}{\ensuremath{\mathrm{FCFF}}}
\newcommand{\FCFE}{\ensuremath{\mathrm{FCFE}}}
\newcommand{\WACC}{\ensuremath{\mathrm{WACC}}}
\newcommand{\NOPAT}{\ensuremath{\mathrm{NOPAT}}}
\newcommand{\ROIC}{\ensuremath{\mathrm{ROIC}}}
\newcommand{\ROE}{\ensuremath{\mathrm{ROE}}}
\newcommand{\ERP}{\ensuremath{\mathrm{ERP}}}
\newcommand{\CRP}{\ensuremath{\mathrm{CRP}}}
\newcommand{\EV}{\ensuremath{\mathrm{EV}}}
\newcommand{\IC}{\ensuremath{\mathrm{IC}}}
\newcommand{\TV}{\ensuremath{\mathrm{TV}}}
\newcommand{\APV}{\ensuremath{\mathrm{APV}}}
\newcommand{\EVA}{\ensuremath{\mathrm{EVA}}}
\newcommand{\Ke}{K_e}
\newcommand{\Kd}{K_d}
\newcommand{\Ku}{K_u}
\newcommand{\Rf}{R_f}
\newcommand{\RR}{\mathrm{RR}}

%% ── Metadados ──────────────────────────────────────────────────────────────────
\title{%
  \textbf{ratvalue} \\[0.5em]
  \large Manual do Usuário \\[0.3em]
  \normalsize Framework de Valuation Damodaran em C++23
}
\author{%
  LAFI --- Laboratório de Avaliação de Firmas e Investimentos \\
  UFPel --- Universidade Federal de Pelotas
}
\date{Junho de 2026 \quad \texttt{v1.0}}

%% ════════════════════════════════════════════════════════════════════════════════
\begin{document}

\maketitle
\tableofcontents

%% ════════════════════════════════════════════════════════════════════════════════
\chapter*{Resumo}
\addcontentsline{toc}{chapter}{Resumo}

\textbf{ratvalue} é uma biblioteca C++23 que implementa o framework completo de
valuation de Aswath Damodaran, desde a normalização de resultados até modelos
especializados para startups, firmas em dificuldade e instituições financeiras.
O projeto é construído sobre \texttt{ratmoney}, que garante aritmética monetária
com precisão racional exata (\texttt{int64\_t} + GCD sobre \texttt{\_\_int128}).

O manual cobre:
\begin{itemize}
  \item Fundamentos teóricos de cada modelo de valuation;
  \item Especificação completa da API em C++23;
  \item Exemplos numéricos verificáveis;
  \item Exemplo real com Petrobras (PETR4) e Itaú Unibanco (ITUB4).
\end{itemize}

Pré-requisitos do leitor: familiaridade com C++17+, álgebra linear básica e
conceitos de finanças corporativas ao nível de um curso de graduação.

%% ════════════════════════════════════════════════════════════════════════════════
\chapter{Introdução}

\section{Motivação}

Modelos de valuation são o núcleo analítico de toda decisão de alocação de
capital.  O framework de Damodaran, desenvolvido ao longo de três décadas em
\textit{Investment Valuation} (3ª ed., 2012) e \textit{The Dark Side of
Valuation} (3ª ed., 2018), é a referência mais rigorosa e completa disponível.

Implementações existentes concentram-se em Python/Excel, que têm dois problemas
para análise quantitativa séria:

\begin{enumerate}
  \item \textbf{Precisão numérica:} \texttt{float64} não representa frações
        decimais exatamente. O erro acumula em cálculos encadeados.
  \item \textbf{Performance:} Python possui GIL e overhead de boxing; Excel não
        é versionável nem auditável como código.
\end{enumerate}

\textbf{ratvalue} resolve ambos: aritmética monetária exata via
\texttt{ratmoney}, e performance de C++ compilado com \texttt{-O2}.

\section{Escopo}

A tabela~\ref{tab:modulos} lista todos os módulos implementados.

\begin{table}[h]
\centering
\caption{Módulos do ratvalue}
\label{tab:modulos}
\begin{tabularx}{\textwidth}{llX}
\toprule
\textbf{Arquivo} & \textbf{Categoria} & \textbf{Conteúdo} \\
\midrule
\texttt{normalize} & Dados & Normalização EBIT (histórica / margem) \\
\texttt{capm}      & Custo de capital & CAPM com CRP (mercados emergentes) \\
\texttt{beta}      & Custo de capital & Hamada; beta bottom-up \\
\texttt{kd}        & Custo de capital & Kd sintético via ICR → rating \\
\texttt{wacc}      & Custo de capital & WACC \\
\texttt{capital\_structure} & Estrutura & Estrutura ótima de capital \\
\texttt{fcff}      & Fluxos & FCFF; projeção multi-estágio \\
\texttt{fcfe}      & Fluxos & FCFE; DCF por equity \\
\texttt{growth}    & Crescimento & ROIC, RR, crescimento fundamental \\
\texttt{dcf}       & Valuation & DCF (FCFF → WACC → EV) \\
\texttt{ddm}       & Valuation & DDM multi-estágio; H-Model \\
\texttt{terminal}  & Valuation & TV consistente; TV por múltiplo \\
\texttt{apv}       & Valuation & APV (MM e Miles-Ezzell) \\
\texttt{excess\_returns} & Análise & EVA / Excess Returns \\
\texttt{relative}  & Análise & P/L, P/VP, EV/EBITDA justificados \\
\texttt{startup}   & Especializados & Valuation por cenários \\
\texttt{distress}  & Especializados & Equity como opção (Merton) \\
\texttt{bank}      & Especializados & FCFE para bancos (capital regulatório) \\
\bottomrule
\end{tabularx}
\end{table}

\section{Dependências e Build}

\textbf{Dependências:}
\begin{itemize}
  \item \texttt{ratmoney} $\geq$ 0.1 (instalado em \texttt{/usr/local})
  \item GCC $\geq$ 13 ou Clang $\geq$ 17 (C++23)
  \item CMake $\geq$ 3.20
  \item GTest (para os testes)
\end{itemize}

\begin{lstlisting}[style=cpp, language=bash, numbers=none,
                   caption={Compilação e execução dos testes}]
cmake -S . -B build -DCMAKE_BUILD_TYPE=Release
cmake --build build -j$(nproc)
ctest --test-dir build --output-on-failure
\end{lstlisting}

%% ════════════════════════════════════════════════════════════════════════════════
\chapter{Infraestrutura: ratmoney e tipos internos}

\section{ratmoney: aritmética monetária exata}

Toda unidade monetária em \textbf{ratvalue} é representada pelo tipo
\texttt{ratmoney::Currency}, que armazena:

\begin{itemize}
  \item \texttt{units}: contagem em menor denominação (\emph{centavos},
        para BRL e USD com \texttt{precision=2});
  \item \texttt{rate}: valor de 1 unidade desta moeda em unidades de uma
        moeda-base implícita (tipo \texttt{Rational});
  \item \texttt{description}: nome, símbolo e precisão decimal.
\end{itemize}

O tipo \texttt{ratmoney::Rational} representa um número racional $p/q$ na
forma canônica: $q > 0$, $\gcd(|p|, q) = 1$.  Toda aritmética usa
\texttt{\_\_int128} como intermediário, eliminando overflow silencioso.

\begin{lstlisting}[style=cpp, caption={Criação de moedas no ratvalue}]
ratmoney::CurrencyDescription BRL{"BRL", "R$", 2};
ratmoney::Rational UNIT_RATE{1, 1};

// R$1 bilhão = 1e11 centavos
Currency B(int64_t bilhoes) {
    return Currency{bilhoes * 100'000'000'000LL, UNIT_RATE, BRL};
}

Currency ebit = B(145);   // R$145 bilhoes
\end{lstlisting}

\begin{atencao}
A identidade de denominação é determinada por \texttt{CurrencyDescription}.
Dois objetos \texttt{Currency} só podem ser somados ou comparados se tiverem
a \emph{mesma descrição}.  O campo \texttt{rate} serve para \emph{conversão
entre moedas}, não para identificar a denominação.
ratvalue valida consistência de moeda em todas as funções relevantes e retorna
\texttt{ValuationError::CurrencyMismatch} em caso de violação.
\end{atencao}

\section{types.h: erros e projeção}

\begin{lstlisting}[style=cpp, caption={Tipos fundamentais do ratvalue}]
enum class ValuationError {
    InvalidInput,     // entrada inválida (negativa, zero onde proibido...)
    WACCLEGrowth,     // WACC <= g: denominador de Gordon seria <= 0
    Overflow,         // overflow aritmético em Rational
    MoneyError,       // CurrencyError do ratmoney
    CurrencyMismatch, // inputs monetários em denominações diferentes
};

struct ProjectionStage {
    ratmoney::Rational growth_rate;
    int                periods;   // anos neste estágio
};
\end{lstlisting}

Todas as funções do ratvalue retornam \texttt{std::expected<T, ValuationError>},
forçando tratamento explícito de erros em tempo de compilação.

\section{detail.h: aritmética racional interna}

O cabeçalho interno \texttt{detail.h} expõe helpers que implementam as
operações $+$, $-$, $\times$, $\div$ sobre \texttt{Rational} com
intermediários \texttt{\_\_int128}:

\begin{lstlisting}[style=cpp]
namespace ratvalue::detail {
    expected<Rational, ValuationError> r_add(Rational a, Rational b);
    expected<Rational, ValuationError> r_sub(Rational a, Rational b);
    expected<Rational, ValuationError> r_mul(Rational a, Rational b);
    expected<Rational, ValuationError> r_div(Rational a, Rational b);
    expected<Rational, ValuationError> one_plus(Rational r); // 1 + r
    expected<Rational, ValuationError> make_rational(__int128 n, __int128 d);

    bool   same_currency(const Currency& a, const Currency& b);
    double to_double(const Currency& c);   // unidades maiores (ex: reais)
    expected<Currency, ValuationError>
        make_currency_from_double(double v, Rational rate,
                                  const CurrencyDescription& desc);
}
\end{lstlisting}

\begin{atencao}
\texttt{detail.h} é um cabeçalho \emph{interno}, não instalado.
Código cliente não deve incluí-lo diretamente.
\end{atencao}

%% ════════════════════════════════════════════════════════════════════════════════
\chapter{Normalização de Resultados}

\section{Teoria}

Firmas cíclicas (petróleo, mineração, siderurgia, agro) apresentam EBIT
volátil que distorce qualquer DCF baseado no ano mais recente.  Damodaran
(cap.~8) recomenda normalizar antes de avaliar, usando uma das duas abordagens:

\begin{description}
  \item[Média histórica]
    Calcula a média aritmética do EBIT dos últimos $N$ anos ($N \geq 2$),
    capturando o ``EBIT do ciclo médio'':
    \[
      EBIT_{\text{norm}} = \frac{1}{N} \sum_{i=1}^{N} EBIT_i
    \]

  \item[Margem normalizada]
    Aplica uma margem operacional histórica (ou setorial) sobre a receita
    atual.  Útil quando a receita é estável mas a margem oscila:
    \[
      EBIT_{\text{norm}} = \text{Receita atual} \times \text{Margem normalizada}
    \]
\end{description}

\section{API}

\begin{lstlisting}[style=cpp]
enum class NormalizationMethod {
    HistoricalAverage,  // método 1
    NormalizedMargin,   // método 2
};

struct NormalizationInputs {
    NormalizationMethod             method;
    std::vector<ratmoney::Currency> historical_ebit;   // >= 2 para Average
    ratmoney::Currency              current_revenue;    // para Margin
    ratmoney::Rational              normalized_margin;  // para Margin
};

expected<Currency, ValuationError>
normalize_ebit(const NormalizationInputs& inputs);
\end{lstlisting}

\begin{exemplo}{Petrobras --- média histórica do ciclo}
\begin{lstlisting}[style=cpp, numbers=none]
auto result = normalize_ebit(NormalizationInputs{
    .method = NormalizationMethod::HistoricalAverage,
    .historical_ebit = {B(80), B(95), B(145), B(165), B(145)},
    // EBIT 2019–2023 em R$ bilhões
    .current_revenue = B(500),  // nao usado neste metodo
    .normalized_margin = {0, 1},
});
// result->units() / 1e11 = R$126B (média)
\end{lstlisting}
\end{exemplo}

\section{Observações}

A aritmética usa \texttt{Currency::add} acumulado e \texttt{Currency::scale}
com fator $\{1, N\}$.  O arredondamento é \textit{HalfEven} (padrão bancário),
de modo que a soma exata é preservada e o erro de arredondamento é de no máximo
1 centavo.

%% ════════════════════════════════════════════════════════════════════════════════
\chapter{Custo de Capital}

\section{CAPM com Risco-País}
\label{sec:capm}

\subsection{Modelo}

Em mercados emergentes, Damodaran (\textit{Investment Valuation}, cap.~7)
estende o CAPM adicionando um prêmio de risco-país ponderado pela exposição
$\lambda$ da empresa:

\begin{equation}
  \Ke = \Rf + \beta \cdot \ERP_{\text{maduro}} + \lambda \cdot \CRP
  \label{eq:capm_crp}
\end{equation}

Onde:
\begin{itemize}
  \item $\Rf$: taxa livre de risco global (\emph{T-Bond 10Y} americano,
        expurgado do componente de inflação diferencial se necessário);
  \item $\ERP_{\text{maduro}}$: prêmio de risco do mercado de referência
        (geralmente EUA; Damodaran publica estimativas anuais);
  \item $\CRP$: prêmio de risco-país, calculado como spread soberano
        $\times$ ($\sigma_{\text{ações}} / \sigma_{\text{bônus}}$);
  \item $\lambda$: exposição da empresa ao risco-país.
        $\lambda = 1$ para empresas com receitas inteiramente domésticas;
        $\lambda < 1$ para exportadoras ou multinacionais.
\end{itemize}

\begin{definicao}{Country Risk Premium (CRP)}
  \[
    \CRP = \text{Spread soberano} \times
           \frac{\sigma_{\text{ações mercado emergente}}}
                {\sigma_{\text{bônus soberano}}}
  \]
  Para o Brasil em 2024: spread EMBI+ $\approx 2\%$;
  fator de volatilidade $\approx 1{,}5$; $\CRP \approx 3\%$.
\end{definicao}

\subsection{API}

\begin{lstlisting}[style=cpp]
struct CAPMInputs {
    Rational risk_free_rate;
    Rational beta;
    Rational equity_risk_premium;              // ERP mercado maduro
    Rational country_risk_premium{0, 1};       // CRP (padrão: 0)
    Rational lambda{1, 1};                     // exposição país (padrão: 1)
};

expected<Rational, ValuationError>
cost_of_equity(const CAPMInputs& inputs);
\end{lstlisting}

Todo o cálculo é feito em aritmética racional exata.  Para
$\Rf = 5\% = \frac{1}{20}$, $\beta = 1{,}30 = \frac{13}{10}$,
$\ERP = 7\% = \frac{7}{100}$, $\CRP = 3\% = \frac{3}{100}$, $\lambda = 1$:

\[
  \Ke
  = \tfrac{1}{20} + \tfrac{13}{10} \cdot \tfrac{7}{100} + 1 \cdot \tfrac{3}{100}
  = \tfrac{50}{1000} + \tfrac{91}{1000} + \tfrac{30}{1000}
  = \tfrac{171}{1000} = 17{,}1\%
\]

\section{Beta: Hamada e Bottom-Up}

\subsection{Equação de Hamada}

A alavancagem financeira amplifica o risco sistemático:

\begin{equation}
  \beta_L = \beta_U \cdot \left[1 + (1-t) \cdot \frac{D}{E}\right]
  \label{eq:hamada}
\end{equation}

Inversamente:
\begin{equation}
  \beta_U = \frac{\beta_L}{1 + (1-t) \cdot D/E}
  \label{eq:hamada_inv}
\end{equation}

\subsection{Beta Bottom-Up}

O beta histórico (regressão de retornos) é ruidoso e reflete a estrutura de
capital passada.  Damodaran recomenda o \emph{beta bottom-up}:

\begin{enumerate}
  \item Coletar $\beta_L$ de $n$ empresas comparáveis (mesmo setor);
  \item Desalavancar cada comparável via eq.~\eqref{eq:hamada_inv};
  \item Calcular a média: $\bar\beta_U = n^{-1}\sum_i \beta_{U,i}$;
  \item Re-alavancar para a estrutura-alvo via eq.~\eqref{eq:hamada}.
\end{enumerate}

\subsection{API}

\begin{lstlisting}[style=cpp]
// Alavancar
expected<Rational, ValuationError>
lever_beta(Rational unlevered_beta, Rational debt_to_equity, Rational tax_rate);

// Desalavancar
expected<Rational, ValuationError>
unlever_beta(Rational levered_beta, Rational debt_to_equity, Rational tax_rate);

struct BottomUpBetaInputs {
    vector<Rational> comparable_levered_betas;
    vector<Rational> comparable_debt_to_equity;
    Rational         comparable_tax_rate;
    Rational         target_debt_to_equity;
    Rational         target_tax_rate;
};

expected<Rational, ValuationError>
bottom_up_beta(const BottomUpBetaInputs& inputs);
\end{lstlisting}

\begin{atencao}
A média dos $\beta_U$ em \texttt{bottom\_up\_beta} é calculada em
\texttt{double}, pois betas de mercado são estimativas sem precisão exata.
A re-alavancagem final retorna \texttt{Rational} com 4 casas decimais.
\end{atencao}

\section{Kd Sintético}
\label{sec:kd_sint}

\subsection{Metodologia}

Empresas sem bônus negociados não têm Kd observável.  Damodaran propõe
estimá-lo a partir do \emph{índice de cobertura de juros} (ICR):

\[
  \text{ICR} = \frac{\text{EBIT}}{\text{Despesa Financeira}}
  \quad\longrightarrow\quad
  \text{rating sintético}
  \quad\longrightarrow\quad
  \Kd = \Rf + \text{spread}(\text{rating})
\]

A tabela de ICR $\to$ rating é publicada anualmente em
\url{https://pages.stern.nyu.edu/~adamodar/}.
O ratvalue implementa as tabelas de 2023 para \emph{large cap} e
\emph{small cap} (thresholds mais exigentes para empresas menores).

\subsection{Tabela resumida (large cap)}

\begin{table}[h]
\centering
\caption{ICR → Rating → Spread (large cap, Damodaran 2023)}
\label{tab:rating}
\begin{tabular}{ccr}
\toprule
\textbf{ICR mínimo} & \textbf{Rating} & \textbf{Spread (bps)} \\
\midrule
$> 8{,}50$ & Aaa/AAA & 63 \\
$6{,}50$   & Aa2/AA  & 78 \\
$5{,}50$   & A1/A+   & 98 \\
$4{,}25$   & A2/A    & 108 \\
$3{,}00$   & A3/A--  & 122 \\
$2{,}50$   & Baa2/BBB & 156 \\
$2{,}25$   & Ba1/BB+ & 200 \\
$2{,}00$   & Ba2/BB  & 240 \\
$1{,}75$   & B1/B+   & 350 \\
$1{,}50$   & B2/B    & 375 \\
$1{,}25$   & B3/B--  & 500 \\
$0{,}80$   & Caa/CCC & 600 \\
$0{,}65$   & Ca2/CC  & 700 \\
$0{,}20$   & C2/C    & 800 \\
$< 0{,}20$ & D2/D    & 1200 \\
\bottomrule
\end{tabular}
\end{table}

\subsection{API}

\begin{lstlisting}[style=cpp]
enum class CreditRating : uint8_t {
    AAA, AA, A_plus, A, A_minus, BBB, BB_plus, BB,
    B_plus, B, B_minus, CCC, CC, C, D
};

struct SyntheticKdInputs {
    Currency ebit;
    Currency interest_expense;
    Rational risk_free_rate;
    bool     large_firm{true};
};

struct SyntheticKdResult {
    Rational     cost_of_debt;      // Kd = Rf + spread (Rational exato)
    Rational     default_spread;
    CreditRating rating;
    double       interest_coverage;
};

expected<SyntheticKdResult, ValuationError>
synthetic_cost_of_debt(const SyntheticKdInputs& inputs);

// versão direta (para o otimizador interno)
expected<SyntheticKdResult, ValuationError>
synthetic_cost_of_debt_from_icr(double icr, Rational rf, bool large_firm);
\end{lstlisting}

\section{WACC}

\subsection{Fórmula}

\begin{equation}
  \WACC = \frac{E}{V} \Ke + \frac{D}{V} \Kd (1-t)
\end{equation}

Onde $V = E + D$.

\subsection{API}

\begin{lstlisting}[style=cpp]
struct WACCInputs {
    Rational equity_weight;   // E/V
    Rational debt_weight;     // D/V
    Rational cost_of_equity;  // Ke
    Rational cost_of_debt;    // Kd
    Rational tax_rate;        // t
};

expected<Rational, ValuationError>
compute_wacc(const WACCInputs& inputs);
\end{lstlisting}

Toda a aritmética é racional exata.  Para os parâmetros do Exemplo Petrobras:
$E_w = 2/3$, $\Ke = 4629/40000$, $D_w = 1/3$, $\Kd = 578/10000$, $t = 17/50$:
\[
  \WACC = \frac{2}{3} \cdot \frac{4629}{40000}
        + \frac{1}{3} \cdot \frac{578}{10000} \cdot \frac{33}{50}
        = \frac{44933}{500000} \approx 8{,}99\%
\]

%% ════════════════════════════════════════════════════════════════════════════════
\chapter{Fluxo de Caixa Livre}

\section{FCFF --- Free Cash Flow to Firm}

\subsection{Fórmula}

\begin{equation}
  \FCFF = \EBIT(1-t) + D\&A - \text{CapEx} - \Delta CGN
  \label{eq:fcff}
\end{equation}

Onde $\Delta CGN$ é a variação no capital de giro não caixa (Damodaran usa
$\Delta \text{NWC} = \Delta \text{CG}$).

\subsection{API}

\begin{lstlisting}[style=cpp]
struct FCFFInputs {
    Currency ebit;
    Rational tax_rate;
    Currency depreciation_amortization;
    Currency capex;
    Currency delta_nwc;
};

expected<Currency, ValuationError>
compute_fcff(const FCFFInputs& inputs);

struct FCFFProjection {
    Currency                    base_fcff;
    std::vector<ProjectionStage> stages;
};

// Retorna [FCFF_1, FCFF_2, ..., FCFF_N] (base não incluído)
expected<std::vector<Currency>, ValuationError>
project_fcff(const FCFFProjection& proj);
\end{lstlisting}

A projeção é exata: cada período aplica \texttt{Currency::scale} com fator
$(1+g)$ representado como \texttt{Rational}.

\section{FCFE --- Free Cash Flow to Equity}

\subsection{Fórmula}

\begin{equation}
  \FCFE = NI - (1-\delta)(\text{CapEx} - D\&A + \Delta CGN) + \Delta D
  \label{eq:fcfe}
\end{equation}

Onde $\delta = D/(D+E)$ é o \emph{debt ratio}.  A interpretação:
o acionista financia apenas a parcela $(1-\delta)$ do reinvestimento líquido;
o restante é coberto por nova dívida.

\subsection{Diferença conceptual: FCFF vs. FCFE}

\begin{table}[h]
\centering
\caption{Comparação FCFF e FCFE}
\begin{tabular}{lll}
\toprule
\textbf{Dimensão} & \textbf{FCFF} & \textbf{FCFE} \\
\midrule
Base & EBIT (pré-juros) & Lucro Líquido (pós-juros) \\
Desconto & WACC & $\Ke$ \\
Resultado & Enterprise Value & Equity Value (direto) \\
Alavancagem & Capturada no WACC & Capturada no FCFE \\
\bottomrule
\end{tabular}
\end{table}

\begin{atencao}
FCFF e FCFE produzem o mesmo valor de equity quando os inputs são consistentes
(mesma estrutura de capital implícita).  Discrepâncias indicam inconsistência
nas hipóteses, não erros no código.
\end{atencao}

\subsection{API}

\begin{lstlisting}[style=cpp]
struct FCFEInputs {
    Currency net_income;
    Rational debt_ratio;                  // delta = D/(D+E)
    Currency capex;
    Currency depreciation_amortization;
    Currency delta_nwc;
    Currency net_new_debt;               // pode ser negativo
};

expected<Currency, ValuationError>
compute_fcfe(const FCFEInputs& inputs);

struct FCFEDCFInputs {
    std::vector<Currency> projected_fcfe;
    Rational              cost_of_equity;
    Rational              terminal_growth;
    int64_t               shares_outstanding;
};

struct FCFEDCFResult {
    Currency equity_value;
    Currency equity_value_per_share;
    Currency pv_fcfe;
    Currency terminal_value;
    Currency pv_terminal_value;
};

expected<FCFEDCFResult, ValuationError>
compute_fcfe_dcf(const FCFEDCFInputs& inputs);
\end{lstlisting}

\section{Crescimento Fundamental}
\label{sec:growth}

\subsection{Teoria}

Damodaran insiste que hipóteses de crescimento devem ser \emph{ancoradas} em
fundamentos.  Crescimento ``gratuito'' (sem reinvestimento) é um erro comum
que infla artificialmente o valor.

\begin{equation}
  g_{\text{firma}} = \ROIC \times \RR
  \label{eq:g_firma}
\end{equation}

\begin{equation}
  g_{\text{equity}} = \ROE \times (1 - \text{payout})
  \label{eq:g_equity}
\end{equation}

Onde:
\[
  \ROIC = \frac{\NOPAT}{\IC} \qquad
  \RR   = \frac{\text{CapEx} - D\&A + \Delta CGN}{\NOPAT}
\]

\subsection{API}

\begin{lstlisting}[style=cpp]
// ROIC = NOPAT / Invested Capital
struct ROICInputs { Currency nopat; Currency invested_capital; };
expected<Rational, ValuationError> compute_roic(const ROICInputs&);

// RR = (CapEx - DA + dNWC) / NOPAT
struct ReinvestmentRateInputs {
    Currency capex, depreciation_amortization, delta_nwc, nopat;
};
expected<Rational, ValuationError>
compute_reinvestment_rate(const ReinvestmentRateInputs&);

// g = ROIC x RR
expected<Rational, ValuationError>
fundamental_growth_firm(Rational roic, Rational reinvestment_rate);

// g = ROE x retention
expected<Rational, ValuationError>
fundamental_growth_equity(Rational roe, Rational retention_ratio);
\end{lstlisting}

%% ════════════════════════════════════════════════════════════════════════════════
\chapter{DCF --- Desconto de Fluxos de Caixa}

\section{Modelo Padrão: FCFF e WACC}

\subsection{Estrutura multi-estágio}

O DCF padrão de Damodaran divide o horizonte em dois estágios:

\begin{enumerate}
  \item \textbf{Período explícito} ($t = 1, \ldots, N$):
        FCFFs projetados descontados individualmente.
  \item \textbf{Perpetuidade estável} ($t > N$):
        valor terminal (TV) calculado via Gordon Growth.
\end{enumerate}

\begin{align}
  \EV
  &= \sum_{t=1}^{N} \frac{\FCFF_t}{(1+\WACC)^t}
   + \frac{\TV_N}{(1+\WACC)^N}
  \label{eq:dcf_ev}
  \\[4pt]
  \TV_N
  &= \frac{\FCFF_N \cdot (1 + g)}{(\WACC - g)}
  \label{eq:gordon}
  \\[4pt]
  \text{Equity} &= \EV - \text{Dívida Líquida}
  \\[4pt]
  P_{\text{ação}} &= \frac{\text{Equity}}{n_{\text{ações}}}
\end{align}

\subsection{Implementação: double para desconto, Rational para TV}

\begin{atencao}
O fator de desconto $(1+r)^{-t}$ é irracional para $t > 1$ na maioria dos
casos.  O ratvalue usa \texttt{double} \emph{apenas} para as potências; todas
as demais operações são racionais exatas.  Em particular, o multiplicador do
TV é um \texttt{Rational} exato:
\[
  \TV_{\text{mul}} = \frac{(1+g)/g_{\text{den}}}{(\WACC - g)}
  = \frac{(g_{\text{den}} + g_{\text{num}}) \cdot \WACC_{\text{den}}}
         {\WACC_{\text{num}} \cdot g_{\text{den}} - g_{\text{num}} \cdot \WACC_{\text{den}}}
\]
\end{atencao}

\subsection{API}

\begin{lstlisting}[style=cpp]
struct DCFInputs {
    std::vector<Currency> projected_fcff;
    Rational              wacc;
    Rational              terminal_growth;
    Currency              net_debt;
    int64_t               shares_outstanding;
};

struct DCFResult {
    Currency enterprise_value;
    Currency equity_value;
    Currency equity_value_per_share;
    Currency pv_fcff;
    Currency terminal_value;      // TV sem descontar (Rational exato)
    Currency pv_terminal_value;
};

expected<DCFResult, ValuationError>
compute_dcf(const DCFInputs& inputs);
\end{lstlisting}

\section{DDM --- Dividend Discount Model}

O DDM substitui FCFF por dividendos e desconta a $\Ke$:

\begin{lstlisting}[style=cpp]
struct DDMInputs {
    Currency                    base_dividend_per_share;
    std::vector<ProjectionStage> stages;
    Rational                    cost_of_equity;
    Rational                    stable_growth;
};

expected<DDMResult, ValuationError>
compute_ddm(const DDMInputs& inputs);
\end{lstlisting}

\section{H-Model (Fuller \& Hsia, 1984)}

O H-Model é uma fórmula fechada para crescimento que \emph{declina
linearmente} de $g_H$ para $g_L$ ao longo de $2H$ anos:

\begin{equation}
  P = \frac{D_0 \cdot \left[(1+g_L) + H \cdot (g_H - g_L)\right]}{\Ke - g_L}
  \label{eq:hmodel}
\end{equation}

Todo o cálculo é \emph{exato} em aritmética racional --- não há desconto
temporal envolvido.

\begin{lstlisting}[style=cpp]
struct HModelInputs {
    Currency base_dividend;
    Rational high_growth;     // g_H
    Rational stable_growth;   // g_L
    Rational half_life;       // H em anos (ex: {5,1})
    Rational cost_of_equity;
};

struct HModelResult {
    Currency intrinsic_value;
    Rational tv_multiplier;   // para auditoria
};

expected<HModelResult, ValuationError>
compute_h_model(const HModelInputs& inputs);
\end{lstlisting}

%% ════════════════════════════════════════════════════════════════════════════════
\chapter{Valor Terminal}

\section{Consistência com ROIC}

\subsection{O problema do Gordon padrão}

O TV pelo Gordon Growth implica uma taxa de reinvestimento:
\[
  \RR_{\text{impl}} = \frac{g}{\ROIC_{\text{estável}}}
\]

Se o analista assume $g = 3\%$ mas $\ROIC_{\text{estável}} = 8\%$, então
$\RR_{\text{impl}} = 37{,}5\%$.  Se ele projetou os FCFFs com $\RR = 30\%$,
há uma descontinuidade no terminal que infla o TV.

\subsection{TV consistente (Damodaran cap.~12)}

\begin{equation}
  \TV = \frac{\NOPAT_{\text{estável}} \cdot (1 - g/\ROIC_{\text{estável}})}{
              \WACC - g}
  \label{eq:tv_consistente}
\end{equation}

Usa o NOPAT do período terminal como base, derivando o FCFF estável pela RR
consistente com o ROIC assumido.

\subsection{Auditoria de consistência}

\begin{lstlisting}[style=cpp]
struct TerminalConsistency {
    Rational implied_reinvestment_rate;  // g / ROIC_estavel
    Rational return_spread;              // ROIC_estavel - WACC
    bool     value_creating;             // ROIC > WACC
    bool     reinvestment_feasible;      // 0 <= g/ROIC <= 1
};

TerminalConsistency check_terminal_consistency(
    Rational wacc, Rational terminal_growth, Rational terminal_roic);

struct ConsistentTVInputs {
    Currency stable_nopat;
    Rational wacc;
    Rational terminal_growth;
    Rational terminal_roic;
};

expected<Currency, ValuationError>
consistent_terminal_value(const ConsistentTVInputs& inputs);
\end{lstlisting}

\begin{atencao}
Quando $\ROIC_{\text{estável}} < \WACC$, o TV ainda é positivo mas a firma
destrói valor na perpetuidade (\texttt{value\_creating = false}).
Isso é economicamente válido; o modelo não rejeita esse caso.
\end{atencao}

\section{TV por Múltiplo}

Alternativa ao Gordon Growth, âncora o TV em um múltiplo de mercado observado:

\begin{equation}
  \TV = \text{EBITDA}_{N} \times \text{EV/EBITDA}_{\text{setor}}
\end{equation}

\begin{lstlisting}[style=cpp]
expected<Currency, ValuationError>
terminal_value_by_multiple(Currency terminal_ebitda,
                            Rational ev_ebitda_multiple);
\end{lstlisting}

%% ════════════════════════════════════════════════════════════════════════════════
\chapter{APV --- Adjusted Present Value}

\section{Motivação}

O \APV\ separa explicitamente o valor operacional do benefício fiscal:

\begin{equation}
  \APV = V_U + PV(\text{Tax Shields})
\end{equation}

Onde $V_U$ é o valor da firma \emph{sem alavancagem}, obtido descontando
os FCFFs ao custo do equity desalavancado $\Ku$:

\[
  \Ku = \Rf + \beta_U \cdot \ERP + \lambda \cdot \CRP
\]

O \APV\ é preferível ao WACC quando a estrutura de capital varia ao longo
do tempo (LBOs, recapitalizações, projetos com cronograma de amortização).

\section{Modigliani-Miller (1963): dívida permanente}

Para dívida permanente e tax shields com risco zero (desconto a $\Rf$):

\begin{equation}
  PV(\text{TS})_{\text{MM}} = t \cdot D
  \label{eq:ts_mm}
\end{equation}

\begin{lstlisting}[style=cpp]
struct APVInputs {
    std::vector<Currency> projected_fcff;
    Rational              terminal_growth;
    Rational              unlevered_cost_of_equity;  // Ku
    Currency              debt_market_value;          // D
    Rational              tax_rate;                  // t
    Currency              net_debt;
    int64_t               shares_outstanding;
};

expected<APVResult, ValuationError>
compute_apv(const APVInputs& inputs);
\end{lstlisting}

\section{Miles-Ezzell (1980): dívida rebalanceada}

Quando a firma mantém um $D/V$ constante (rebalanceia a cada período), os
tax shields têm o risco dos ativos e devem ser descontados a $\Ku$, exceto
o primeiro período (conhecido com certeza):

\begin{equation}
  PV(\text{TS})_{\text{ME}}
  = \frac{t \cdot \Kd \cdot D \cdot (1 + \Ku)}{(1 + \Kd)(\Ku - g)}
  \label{eq:ts_me}
\end{equation}

O \textbf{multiplicador de Miles-Ezzell} é calculado em aritmética racional
exata via \texttt{detail::r\_mul} e \texttt{detail::r\_div}.

\begin{lstlisting}[style=cpp]
struct APVMilesEzzellInputs {
    // ... todos os campos de APVInputs, mais:
    Rational cost_of_debt;   // Kd (necessario para formula ME)
};

expected<APVResult, ValuationError>
compute_apv_miles_ezzell(const APVMilesEzzellInputs& inputs);
\end{lstlisting}

\begin{exemplo}{Comparação MM vs. ME (Petrobras)}
$D = \text{R\$}250B$, $t = 34\%$, $\Kd = 5{,}78\%$, $\Ku \approx 11{,}57\%$, $g = 1{,}5\%$:
\[
  PV(\text{TS})_{\text{MM}} = 0{,}34 \times 250 = \text{R\$}85B
\]
\[
  PV(\text{TS})_{\text{ME}}
  = \frac{0{,}34 \times 0{,}0578 \times 250 \times 1{,}1157}{1{,}0578 \times 0{,}1007}
  \approx \text{R\$}51B
\]
O MM superestima o benefício fiscal em $\approx 40\%$ neste caso.
\end{exemplo}

%% ════════════════════════════════════════════════════════════════════════════════
\chapter{EVA e Retornos em Excesso}

\section{Fundamento teórico}

O modelo de \emph{excess returns} (Damodaran, cap.~13) decompõe o valor
da firma em:

\begin{equation}
  V = \IC + PV(\EVA)
  \label{eq:eva_decomp}
\end{equation}

Onde $\EVA_t = (\ROIC_t - \WACC) \times \IC_t$ e $\IC_t$ cresce com o
reinvestimento.  Para o estágio estável (single-stage, Gordon):

\begin{equation}
  V = \IC \cdot \frac{\ROIC - g}{\WACC - g}
  \label{eq:eva_single}
\end{equation}

\textbf{Equivalência com DCF:} quando os inputs são consistentes
($g = \ROIC \times \RR$ e $\FCFF = \NOPAT(1-\RR)$), as equações
\eqref{eq:dcf_ev} e \eqref{eq:eva_decomp} produzem o mesmo $V$.
A decomposição do \EVA\ revela o \emph{valor dos ativos existentes} vs.\
o \emph{valor do crescimento futuro}.

\section{Interpretação econômica}

\begin{itemize}
  \item $\ROIC > \WACC$: a firma cria valor acima do custo de capital;
        $PV(\EVA) > 0$.
  \item $\ROIC = \WACC$: crescimento \emph{neutro}; $V = \IC$.
  \item $\ROIC < \WACC$: cada real investido destrói valor; $PV(\EVA) < 0$.
\end{itemize}

\section{API}

\begin{lstlisting}[style=cpp]
struct SingleStageERInputs {
    Currency invested_capital;
    Rational roic;
    Rational wacc;
    Rational terminal_growth;
};

struct ExcessReturnsResult {
    Currency firm_value;
    Currency asset_value;        // = IC
    Currency pv_excess_returns;  // PV(EVA)
};

expected<ExcessReturnsResult, ValuationError>
excess_returns_value(const SingleStageERInputs& inputs);

struct MultiStageERInputs {
    Currency                    base_invested_capital;
    Rational                    roic;
    std::vector<ProjectionStage> stages;
    Rational                    wacc;
    Rational                    terminal_roic;
    Rational                    terminal_growth;
};

expected<ExcessReturnsResult, ValuationError>
excess_returns_multistage(const MultiStageERInputs& inputs);
\end{lstlisting}

%% ════════════════════════════════════════════════════════════════════════════════
\chapter{Múltiplos Justificados}

\section{Fundamento}

Múltiplos de mercado (P/L, P/VP, EV/EBITDA) são frequentemente interpretados
sem referência aos fundamentos.  Damodaran (cap.~17--20) deriva o
\emph{múltiplo justificado}: o valor que seria correto dado o crescimento,
a rentabilidade e o custo de capital da firma.

\begin{align}
  P/L &= \frac{\text{payout}}{\Ke - g_{\text{equity}}} \label{eq:pe} \\[4pt]
  P/VP &= \ROE \times \frac{\text{payout}}{\Ke - g_{\text{equity}}} \label{eq:pb} \\[4pt]
  EV/EBITDA &= \frac{(1-t)(1-\RR)}{\WACC - g} \label{eq:ev_ebitda}
\end{align}

Onde $g_{\text{equity}} = \ROE \times (1 - \text{payout})$ e
$\RR = g/\ROIC$ (taxa de reinvestimento da firma).

\section{API}

\begin{lstlisting}[style=cpp]
struct JustifiedMultiplesInputs {
    Rational cost_of_equity;
    Rational wacc;
    Rational roe;
    Rational payout_ratio;
    Rational roic;
    Rational tax_rate;
};

struct JustifiedMultiples {
    Rational pe_ratio;
    Rational pb_ratio;
    Rational ev_ebitda;
    Rational g_equity;
    Rational implied_reinvestment;
};

expected<JustifiedMultiples, ValuationError>
compute_justified_multiples(const JustifiedMultiplesInputs& inputs);
\end{lstlisting}

\begin{exemplo}{Petrobras --- múltiplos justificados (ROE normalizado = 15\%)}
Com $\Ke = 11{,}57\%$, $\WACC = 8{,}99\%$, $\ROE_{\text{norm}} = 15\%$,
payout $= 60\%$, $\ROIC = 16{,}63\%$, $t = 34\%$:
\[
  g = 15\% \times 40\% = 6\%
\]
\[
  P/L = \frac{0{,}60}{0{,}1157 - 0{,}06} = 10{,}8\times
  \qquad
  P/VP = 0{,}15 \times 10{,}8 = 1{,}6\times
  \qquad
  EV/EBITDA = 14{,}1\times
\]
\end{exemplo}

%% ════════════════════════════════════════════════════════════════════════════════
\chapter{Estrutura Ótima de Capital}

\section{Algoritmo}

O custo de capital mínimo maximiza o valor da firma.  O ratvalue varre
razões de endividamento $d \in \{0/N, 1/N, \ldots, (N-1)/N\}$ e para
cada ponto:

\begin{enumerate}
  \item Calcula $D/E = d/(1-d)$ como \texttt{Rational};
  \item Alavancar $\beta_U \to \beta_L$ via Hamada (eq.~\ref{eq:hamada});
  \item Calcula $\Ke$ via CAPM (eq.~\ref{eq:capm_crp});
  \item Itera $\Kd$ a partir do ICR:
        \begin{enumerate}
          \item $\text{Interesse} = D \times \Kd$
          \item $\text{ICR} = \text{EBIT} / \text{Interesse}$
          \item $\Kd \leftarrow \Rf + \text{spread}(\text{ICR})$
          \item Repetir até convergência ($< 10^{-9}$, usualmente 2--5 iterações)
        \end{enumerate}
  \item Calcula $\WACC$ exato como \texttt{Rational}.
\end{enumerate}

O ótimo é o ponto de mínimo $\WACC$.

\section{API}

\begin{lstlisting}[style=cpp]
struct OptimalCapitalStructureInputs {
    Rational  unlevered_beta;
    Rational  risk_free_rate;
    Rational  equity_risk_premium;
    Rational  country_risk_premium{0, 1};
    Rational  lambda{1, 1};
    Rational  tax_rate;
    Currency  ebit;       // para calcular ICR
    Currency  firm_value; // E+D (base para D = d * FV)
    int       steps{10};
    bool      large_firm{true};
};

struct CapitalStructurePoint {
    Rational     debt_ratio;
    Rational     levered_beta;
    Rational     cost_of_equity;
    Rational     cost_of_debt;
    CreditRating rating;
    double       interest_coverage;
    Rational     wacc;
};

struct OptimalCapitalStructureResult {
    std::vector<CapitalStructurePoint> schedule;
    CapitalStructurePoint              optimal;
};

expected<OptimalCapitalStructureResult, ValuationError>
optimal_capital_structure(const OptimalCapitalStructureInputs& inputs);
\end{lstlisting}

\begin{atencao}
O modelo de Damodaran via ICR captura apenas o \emph{custo de default
via spread}.  Custos indiretos de dificuldade financeira (perda de
clientes, contratação subótima, distraction gerencial) não são modelados.
Em firmas com EBIT alto relativo ao valor, o algoritmo pode indicar
razões de endividamento elevadas que na prática seriam inviáveis.
\end{atencao}

%% ════════════════════════════════════════════════════════════════════════════════
\chapter{Modelos Especializados}

\section{Valuation por Cenários: Startups e Firmas em Transição}

\subsection{Motivação}

Empresas jovens ou em transição tecnológica têm distribuição de resultados
fortemente assimétrica.  Um único DCF ``base'' ignora essa dispersão.
Damodaran (\textit{Dark Side}, cap.~10--12) recomenda:

\begin{enumerate}
  \item Definir $n$ cenários (geralmente 3--5) com probabilidades;
  \item Calcular o TV de cada cenário: $\TV_i = \FCFF_i \cdot \text{mul}_i$;
  \item Descontar cada TV ao presente: $PV_i = \TV_i / (1+\WACC_i)^T$;
  \item Aplicar probabilidade de sobrevivência:
\end{enumerate}

\begin{equation}
  \EV_{\text{ajustado}} = \left[\sum_i P_i \cdot PV_i \right]
                          \cdot P(\text{sobrevive})
                        + V_{\text{falência}} \cdot [1 - P(\text{sobrevive})]
\end{equation}

\subsection{API}

\begin{lstlisting}[style=cpp]
struct ValuationScenario {
    std::string name;
    Rational    probability;
    Currency    terminal_fcff;
    Rational    terminal_growth;
    Rational    wacc;
    int         years_to_terminal;
};

struct StartupValuationInputs {
    std::vector<ValuationScenario> scenarios;
    Rational  survival_probability;
    Currency  failure_value;
    Currency  net_debt;
    int64_t   shares_outstanding;
};

expected<StartupValuationResult, ValuationError>
startup_valuation(const StartupValuationInputs& inputs);
\end{lstlisting}

A soma das probabilidades deve ser $1{,}0$ (tolerância $10^{-6}$);
caso contrário, \texttt{InvalidInput} é retornado.

\section{Equity como Opção: Modelo de Merton}

\subsection{Fundamento}

Em firmas com dívida elevada, o equity comporta-se como uma \emph{opção de
compra} sobre o valor dos ativos, com strike igual ao valor de face da dívida
(Merton, 1974):

\begin{align}
  d_1 &= \frac{\ln(V/D) + (r + \sigma_V^2/2)T}{\sigma_V \sqrt{T}} \\[4pt]
  d_2 &= d_1 - \sigma_V \sqrt{T} \\[4pt]
  E   &= V \cdot \mathcal{N}(d_1) - D \cdot e^{-rT} \cdot \mathcal{N}(d_2)
  \label{eq:merton}
\end{align}

A probabilidade \emph{risk-neutral} de default é $\mathcal{N}(-d_2)$.
A distância ao default $d_2$ é o indicador de saúde financeira.

\begin{atencao}
$V$ e $\sigma_V$ são os ativos da firma, \emph{não observáveis} diretamente.
A estimação iterativa (resolver $\{E, \sigma_E\}$ com o sistema não-linear
$E = \text{BS}(V, \sigma_V)$ e $\sigma_E E = \mathcal{N}(d_1) \sigma_V V$)
não é implementada; o usuário fornece $V$ e $\sigma_V$.
\end{atencao}

\subsection{API}

\begin{lstlisting}[style=cpp]
struct EquityAsOptionInputs {
    double firm_value;        // V
    double debt_face_value;   // D
    double asset_volatility;  // sigma_V
    double risk_free_rate;    // r (continuo)
    double debt_maturity;     // T em anos
};

struct EquityAsOptionResult {
    double equity_value;
    double debt_market_value;
    double probability_default;  // N(-d2)
    double distance_to_default;  // d2
    double d1;
};

EquityAsOptionResult                         // sem expected: sempre definido
equity_as_option(const EquityAsOptionInputs& inputs) noexcept;
\end{lstlisting}

A CDF normal é implementada via:
$\mathcal{N}(x) = \tfrac{1}{2}\,\mathrm{erfc}\!\left(-x/\sqrt{2}\right)$
(função \texttt{std::erfc} de \texttt{<cmath>}).

\section{FCFE para Instituições Financeiras}

\subsection{Por que bancos são diferentes}

\begin{itemize}
  \item A dívida (depósitos, CDBs) é \emph{matéria-prima}, não fonte de
        financiamento capex --- não faz sentido calcular FCFF ou WACC.
  \item O reinvestimento é regulado: Basileia III exige capital mínimo
        sobre os ativos ponderados pelo risco (RWA).
  \item O custo de capital relevante é apenas $\Ke$.
\end{itemize}

\subsection{Modelo}

\begin{equation}
  \FCFE_{\text{banco}} = NI \times (1 - \RR_{\text{equity}})
\end{equation}

\[
  \RR_{\text{equity}} = \frac{\Delta \text{RWA} \times \text{target Tier1}}{NI}
\]

\subsection{API}

\begin{lstlisting}[style=cpp]
struct BankFCFEInputs {
    Currency net_income;
    Rational equity_reinvestment_rate;
};

expected<Currency, ValuationError>
compute_bank_fcfe(const BankFCFEInputs& inputs);

struct RegulatoryCapitalInputs {
    Currency net_income;
    Currency current_rwa;
    Currency projected_rwa;
    Rational target_tier1_ratio;
};

expected<Rational, ValuationError>
bank_equity_reinvestment_rate(const RegulatoryCapitalInputs& inputs);
\end{lstlisting}

Para o DCF de um banco, usa-se \texttt{compute\_fcfe\_dcf} com os FCFEs
gerados por \texttt{compute\_bank\_fcfe}.

%% ════════════════════════════════════════════════════════════════════════════════
\chapter{Exemplo Completo: Petrobras e Itaú Unibanco}

\section{Dados}

\begin{table}[h]
\centering
\caption{Petrobras S.A.\ --- dados aproximados 2023}
\begin{tabularx}{\textwidth}{Xrr}
\toprule
\textbf{Item} & \textbf{R\$ bilhões} & \textbf{Código} \\
\midrule
Receita líquida & 500 & \texttt{B(500)} \\
EBIT histórico 2019--2023 & 80, 95, 145, 165, 145 & \texttt{B(80)\ldots} \\
EBIT normalizado (média) & 126 & --- \\
D\&A & 45 & \texttt{B(45)} \\
CapEx & 70 & \texttt{B(70)} \\
$\Delta$CGN & 5 & \texttt{B(5)} \\
Lucro Líquido & 124 & \texttt{B(124)} \\
Despesa Financeira & 22 & \texttt{B(22)} \\
Dívida Líquida & 250 & \texttt{B(250)} \\
Capital Investido & 500 & \texttt{B(500)} \\
Ações & 13 bilhões & \texttt{13e9} \\
\bottomrule
\end{tabularx}
\end{table}

\section{Parâmetros de Custo de Capital}

\begin{itemize}
  \item $\Rf = 5\%$ (T-Bond 10Y \emph{stripped});
  \item $\ERP = 5\%$ (mercado americano, Damodaran 2024);
  \item $\CRP = 3\%$ (Brasil, EMBI+ ajustado por volatilidade);
  \item $\lambda = 0{,}5$ (Petrobras: $\approx 50\%$ receita no exterior);
  \item $\beta_U = 0{,}8$ (setor O\&G, bottom-up a partir de Shell, Exxon,
        BP, Chevron, Equinor com $t_{\text{setor}} = 25\%$);
  \item $D/E = 250/500 = 0{,}5$; $t = 34\%$.
\end{itemize}

\textbf{Resultados de custo de capital:}
$\beta_L \approx 1{,}014$; $\Ke \approx 11{,}57\%$;
ICR$= 145/22 = 6{,}6\times \to$ AA rating $\to \Kd = 5{,}78\%$;
$\WACC \approx 8{,}99\%$ (racional exato: $44933/500000$).

\section{Crescimento e FCFF}

\begin{align*}
  \FCFF_0 &= 126 \times 0{,}66 + 45 - 70 - 5 = \text{R\$}53{,}2B \\
  \ROIC   &= 83{,}2/500 = 16{,}63\% \\
  \RR     &= 30/83{,}2 = 36{,}08\% \\
  g_{\text{fund}} &= 16{,}63\% \times 36{,}08\% = 6\%
\end{align*}

Para o DCF usa-se $g_{\text{explícito}} = 3\%$ por 5 anos e
$g_{\text{terminal}} = 1{,}5\%$.

\section{Resumo dos Resultados}

\begin{table}[h]
\centering
\caption{Comparativo de modelos --- Petrobras (PETR4)}
\label{tab:resultados}
\begin{tabular}{lrr}
\toprule
\textbf{Modelo} & \textbf{Preço justo} & \textbf{Módulo} \\
\midrule
DCF (FCFF, Gordon TV)  & R\$39,88 & \texttt{dcf} \\
FCFE (equity direto)   & R\$85,72 & \texttt{fcfe} \\
APV Modigliani-Miller  & R\$31,13 & \texttt{apv} \\
APV Miles-Ezzell       & R\$28,55 & \texttt{apv} \\
H-Model (dividendos)   & R\$54,11 & \texttt{ddm} \\
Cenários (transição E) & R\$24,42 & \texttt{startup} \\
EVA / Excess Returns   & R\$58,51 & \texttt{excess\_returns} \\
\midrule
Cotação mercado (2024) & R\$38--42 & --- \\
\bottomrule
\end{tabular}
\end{table}

\begin{atencao}
A discrepância entre modelos é \emph{esperada e informativa}:
\begin{itemize}
  \item O \textbf{FCFE} usa NI de 2023 (R\$124B) sem normalização --- inflado pelo ciclo
        de petróleo elevado;
  \item O \textbf{APV} não inclui benefício da dívida nas taxas de crescimento,
        de modo que o TV é menor;
  \item O modelo de \textbf{cenários} penaliza fortemente a transição energética
        acelerada (20\% de probabilidade de FCFF cair para R\$40B);
  \item O \textbf{DCF/FCFF normalizado} (R\$39,88) é o mais consistente com
        a cotação de mercado por usar EBIT médio do ciclo.
\end{itemize}
\end{atencao}

\section{Banco Itaú Unibanco (ITUB4)}

\begin{table}[h]
\centering
\caption{Itaú Unibanco --- dados aproximados 2023}
\begin{tabular}{lr}
\toprule
\textbf{Item} & \textbf{Valor} \\
\midrule
Lucro Líquido  & R\$35B \\
RWA atual      & R\$1{.}800B \\
RWA projetado  & R\$1{.}944B (+8\%) \\
Target Tier1   & 13{,}5\% \\
\midrule
$\Delta$Capital regulatório & R\$144B $\times$ 13,5\% = R\$19,4B \\
$\RR_{\text{equity}}$       & R\$19,4B / R\$35B = 55,5\% \\
$\FCFE_{\text{banco}}$      & R\$35B $\times$ 44,5\% = R\$15,6B \\
Preço justo (DCF)           & R\$22,61/ação \\
\bottomrule
\end{tabular}
\end{table}

A alta $\RR_{\text{equity}}$ reflete o crescimento rápido do crédito brasileiro:
o banco precisa reter $\approx 56\%$ dos lucros para suportar a expansão
regulatória dos ativos.

%% ════════════════════════════════════════════════════════════════════════════════
\chapter{Fluxo de Dados entre Módulos}

\section{Visão geral}

A figura abaixo ilustra o fluxo típico de uma análise completa:

\begin{center}
\small
\begin{tabular}{c}
\texttt{normalize} \\
$\downarrow$ \\
\texttt{fcff} $\;\longrightarrow\;$ \texttt{growth}
             $\;\longrightarrow\;$ \texttt{project\_fcff} \\
$\downarrow$ \\
\texttt{beta} + \texttt{capm} + \texttt{kd}
$\;\longrightarrow\;$ \texttt{wacc} \\
$\downarrow$ \\
\texttt{dcf} $\;\longrightarrow\;$ \texttt{terminal} \\
$\quad\longrightarrow\;$ \texttt{apv} (MM + ME) \\
$\quad\longrightarrow\;$ \texttt{excess\_returns} \\
$\quad\longrightarrow\;$ \texttt{relative} \\
$\quad\longrightarrow\;$ \texttt{capital\_structure} \\
\texttt{fcfe} $\;\longrightarrow\;$ \texttt{fcfe\_dcf} \\
\texttt{ddm}  $\;\longrightarrow\;$ \texttt{h\_model} \\
\texttt{startup} (independente) \\
\texttt{distress} (independente) \\
\texttt{bank} $\;\longrightarrow\;$ \texttt{fcfe\_dcf}
\end{tabular}
\end{center}

\section{Consistência entre modelos}

Para que FCFF-DCF e FCFE-DCF produzam o mesmo preço por ação:
\begin{enumerate}
  \item $\FCFE = \FCFF - \text{Juros}(1-t) + \Delta D$;
  \item A taxa de crescimento usada em ambos deve ser a mesma;
  \item O $\Delta D$ deve ser consistente com a razão $D/V$ ao longo do tempo.
\end{enumerate}

Para que APV e WACC-DCF produzam o mesmo $\EV$:
\begin{itemize}
  \item WACC presume alavancagem constante (ajusta WACC a cada período);
  \item APV presume cronograma de dívida predeterminado.
  \item Em geral, APV é mais flexível; WACC é mais simples para estrutura estável.
\end{itemize}

%% ════════════════════════════════════════════════════════════════════════════════
\appendix

\chapter{Tabela de Ratings e Spreads (Damodaran 2023)}

\begin{table}[h]
\centering
\caption{Large Cap}
\begin{tabular}{crr}
\toprule
\textbf{Rating} & \textbf{ICR mínimo} & \textbf{Spread (bps)} \\
\midrule
Aaa/AAA & 8,50 & 63 \\
Aa2/AA  & 6,50 & 78 \\
A1/A+   & 5,50 & 98 \\
A2/A    & 4,25 & 108 \\
A3/A--  & 3,00 & 122 \\
Baa2/BBB & 2,50 & 156 \\
Ba1/BB+ & 2,25 & 200 \\
Ba2/BB  & 2,00 & 240 \\
B1/B+   & 1,75 & 350 \\
B2/B    & 1,50 & 375 \\
B3/B--  & 1,25 & 500 \\
Caa/CCC & 0,80 & 600 \\
Ca2/CC  & 0,65 & 700 \\
C2/C    & 0,20 & 800 \\
D2/D    & $<0{,}20$ & 1200 \\
\bottomrule
\end{tabular}
\quad
\begin{tabular}{crr}
\toprule
\textbf{Rating} & \textbf{ICR mínimo} & \textbf{Spread (bps)} \\
\midrule
Aaa/AAA & 12,50 & 63 \\
Aa2/AA  &  9,50 & 78 \\
A1/A+   &  7,50 & 98 \\
A2/A    &  6,00 & 108 \\
A3/A--  &  4,50 & 122 \\
Baa2/BBB &  4,00 & 156 \\
Ba1/BB+ &  3,50 & 200 \\
Ba2/BB  &  3,00 & 240 \\
B1/B+   &  2,50 & 350 \\
B2/B    &  2,00 & 375 \\
B3/B--  &  1,50 & 500 \\
Caa/CCC &  1,25 & 600 \\
Ca2/CC  &  0,80 & 700 \\
C2/C    &  0,50 & 800 \\
D2/D    & $<0{,}50$ & 1200 \\
\bottomrule
\end{tabular}
\end{table}
\begin{center}\small Esquerda: \emph{large cap}. \quad Direita: \emph{small cap}.\end{center}

\chapter{Referências e Leitura Adicional}

\begin{thebibliography}{99}

\bibitem{dam12}
Damodaran, A. (2012).
\textit{Investment Valuation: Tools and Techniques for Determining the Value of Any Asset}.
3ª ed. Wiley Finance.

\bibitem{dam09}
Damodaran, A. (2009).
\textit{The Dark Side of Valuation}.
2ª ed. FT Press.
[3ª ed., 2018]

\bibitem{me80}
Miles, J.\ A.; Ezzell, J.\ R. (1980).
The Weighted Average Cost of Capital, Perfect Capital Markets, and Project Life:
A Clarification.
\textit{Journal of Financial and Quantitative Analysis}, 15(3), 719--730.

\bibitem{mm63}
Modigliani, F.; Miller, M.\ H. (1963).
Corporate Income Taxes and the Cost of Capital: A Correction.
\textit{American Economic Review}, 53(3), 433--443.

\bibitem{merton74}
Merton, R.\ C. (1974).
On the Pricing of Corporate Debt: The Risk Structure of Interest Rates.
\textit{Journal of Finance}, 29(2), 449--470.

\bibitem{fh84}
Fuller, R.\ J.; Hsia, C.\ C. (1984).
A Simplified Common Stock Valuation Model.
\textit{Financial Analysts Journal}, 40(5), 49--56.

\bibitem{hamada72}
Hamada, R.\ S. (1972).
The Effect of the Firm's Capital Structure on the Systematic Risk of Common Stocks.
\textit{Journal of Finance}, 27(2), 435--452.

\bibitem{dam_site}
Damodaran, A.
\textit{Damodaran Online} --- dados e tabelas atualizados anualmente.
\url{https://pages.stern.nyu.edu/~adamodar/}

\bibitem{ratmoney}
\textit{ratmoney} --- Biblioteca C++23 para aritmética monetária racional exata.
\url{https://github.com/sheep-farm/ratmoney}

\end{thebibliography}

\end{document}
